% ==== P2 draft f1 — §2.2 대체 블록 (L106–L219 대체) ====
\subsection{단일 삽입 자리 --- 점유 통계와 그 요동}\label{sec:sm-site}
앞 소절이 놓은 대정준 도구(식~\eqref{eq:sm-gc})를 가장 작은 계 --- 격자의 삽입 자리 \emph{하나} --- 에
적용한다. 자리는 전해질 너머의 Li 저장조와 입자를 교환하므로 식~\eqref{eq:sm-gc} 의 대정준 설정 그대로다.
전개는 두 단으로 놓는다 --- 먼저 이온을 내부 구조 없는 점입자로 이상화한 두-상태 원형을 표준 유도
그대로 닫고, 이어 점유 이온의 진동 자유도를 켠 확장이 같은 함수형으로 되돌아옴을 보인 뒤, 점유의
요동까지 읽는다. 도착점인 점유율 $\theta$ 와 유효 자리 자유에너지 $\tilde\varepsilon(T)$ 는
\S\ref{sec:sm-lattice} 의 lattice gas 와 \S\ref{sec:sm-electro} 의 전기화학 결선이 받는 입력이다.

\textbf{(a) 출발 --- 두 가정과 원형의 두-상태 목록.} 유도에 앞서 두 가정을 명시한다.
\begin{itemize}[leftmargin=2.2em,itemsep=0.1em]
\item \textbf{가정 1 (hard-core 배타).} 한 자리에는 Li 이온이 $0$ 개 또는 $1$ 개만 들어간다. 근거는 양자
통계가 아니라 강한 단거리 반발이다 --- 자리의 부피와 이온 사이 정전$\cdot$입체 반발이 이중 점유의
에너지를 사실상 무한대로 올린다. 이 가정이 점유수 합을 $n\in\{0,1\}$ 로 자른다.
\item \textbf{가정 2 (자리 독립).} 이 소절에서는 자리의 에너지가 이웃 자리의 점유와 무관하다고 둔다.
이 가정은 \S\ref{sec:sm-mf} 에서 상호작용 파라미터 $\Omega$ 로 완화된다.
\end{itemize}
원형에서는 여기에 이상화 하나를 더 얹는다 --- 이온을 내부 구조 없는 점입자로 두어, 점유 이온이 우물
안에서 갖는 위치$\cdot$운동량 배치를 잠시 얼린다(이 이상화는 아래 확장이 걷어낸다). 그러면 미시상태의
목록은 단 둘이다: 빈 자리($n=0$, 에너지 $0$ 기준)와 점유 자리($n=1$, 에너지 $\varepsilon_0$).

\textbf{(b) 연산 --- 두 항의 대정준 합.} 각 미시상태의 무게는 Gibbs 인자
$e^{-\beta(E_i-\mu N_i)}$(식~\eqref{eq:sm-gibbsfactor})이고, 목록이 둘뿐이라 대정준 합도 두 항이다 ---
빈 자리가 $1$, 점유 자리가 $e^{-\beta(\varepsilon_0-\mu)}$:
\begin{equation}
\Xi_1^{0}\;=\;\sum_{n=0,1}e^{-\beta(\varepsilon_0-\mu)\,n}
\;=\;1+e^{-\beta(\varepsilon_0-\mu)}.
\label{eq:sm-baresum}
\end{equation}
기호에 관하여 --- $\Xi$ 는 자리 1개의 \emph{대정준} 분배함수로, \S\ref{sec:sm-ensemble} 의 canonical
$Z$ 와는 앙상블이 다르므로 기호부터 구분해 적고(첨자 $1$ = 자리 1개), 위 첨자 $0$ 은 내부 자유도를 얼린
원형임을 표시한다.

\textbf{(c) 중간식 --- 상태 확률과 평균 점유.} 식~\eqref{eq:sm-gc} 의 둘째 식으로 두 상태의 확률은
$P_0=1/\Xi_1^{0}$·$P_1=e^{-\beta(\varepsilon_0-\mu)}/\Xi_1^{0}$ 이고, 평균 점유는 그 확률 가중 평균이다:
\begin{equation}
\langle n\rangle
\;=\;0\cdot P_0+1\cdot P_1
\;=\;\frac{e^{-\beta(\varepsilon_0-\mu)}}{1+e^{-\beta(\varepsilon_0-\mu)}},
\label{eq:sm-baremid}
\end{equation}
그리고 식~\eqref{eq:sm-gc} 의 셋째 식 $k_BT\,\partial\ln\Xi_1^{0}/\partial\mu$ 를
식~\eqref{eq:sm-baresum} 에 적용해도 같은 결과가 나온다(교차검증).

\textbf{(d) 박스 --- 두-상태 원형 $\Xi_1^{0}$·$\langle n\rangle$.} 식~\eqref{eq:sm-baremid} 의
분자$\cdot$분모를 $e^{-\beta(\varepsilon_0-\mu)}$ 로 나누면
\begin{equation}
\boxed{\;\Xi_1^{0}\;=\;1+e^{-\beta(\varepsilon_0-\mu)},\qquad
\langle n\rangle\;=\;\frac{1}{1+e^{+\beta(\varepsilon_0-\mu)}}\;.}
\label{eq:sm-bare}
\end{equation}
이 두-상태 대정준 유도는 국소화 흡착(localized adsorption)의 Langmuir 등온선 표준 유도
그대로이고 \cite{hill1960,fowler1939}, 삽입 전극에의 적용 원형은 McKinnon--Haering
\cite{mckinnon1983} 이다. 식~\eqref{eq:sm-bare} 의 $\langle n\rangle$ 은 준위 하나의 평균 전자 점유를
주는 Fermi--Dirac 함수 $f(E)=1/(e^{\beta(E-\mu)}+1)$ 에서 $E$ 자리에 $\varepsilon_0$ 이 앉은 것과
\emph{동형}이다 --- 이 대응의 물리적 배경(페르미온$\cdot$보손의 양자 통계)은 별도의 배경 박스가
자족적으로 정리한다.
\begin{verifybox}
극한 검산 --- (i) $\beta\to\infty$($T\to0$): $\varepsilon_0\gtrless\mu$ 에 따라 $\langle n\rangle\to0/1$
의 계단 --- $E_i-\mu N_i$ 가 낮은 상태만 남는다는 \S\ref{sec:sm-ensemble} 의 검산과 합치한다.
(ii) $\beta\to0$: $\langle n\rangle\to\tfrac12$ --- 두 미시상태 등확률. (iii) $\mu\to\pm\infty$:
$\langle n\rangle\to1/0$ --- 저장조가 이온을 한없이 밀어 넣는/빼내는 극한. 부호 읽기 --- 분모의 지수가
$+\beta(\varepsilon_0-\mu)$ 이므로 점유가 유리할수록($\varepsilon_0<\mu$) 지수가 음으로 죽어
$\langle n\rangle$ 이 $1$ 에 붙는 방향이다.
\end{verifybox}

\begin{bgbox}
\textbf{페르미온·보손과 양자 통계 ---} 양자역학에서 같은 종류의 입자 둘은
\emph{동일입자}(identical particles)다 --- 전자 둘을 맞바꿔도 어떤 측정값도 달라지지 않는다. 이
구별불가능성은 다입자 파동함수에 교환 대칭 제약을 건다: 두 입자를 맞바꿀 때 파동함수가 그대로인
\emph{대칭} 부류와 부호가 뒤집히는 \emph{반대칭} 부류, 자연에는 이 둘만 존재한다. 대칭 쪽이
\emph{보손}(boson)이다 --- 정수 스핀을 가지며 한 양자 상태를 몇 개라도 함께 점유할 수 있어, 준위 $E$ 의
평형 평균 점유수가 Bose--Einstein 분포 $\bar n(E)=1/(e^{\beta(E-\mu)}-1)$ 를 따른다. 반대칭 쪽이
\emph{페르미온}(fermion)이다 --- 반정수 스핀을 가지며 같은 상태에 둘을 넣으면 파동함수가 자기 자신과
상쇄되어 사라지므로 준위당 점유가 $0$ 또는 $1$ 로 제한되고(\emph{Pauli 배타 원리}, Pauli exclusion
principle), 평균 점유가 Fermi--Dirac 분포 $f(E)=1/(e^{\beta(E-\mu)}+1)$ 를 따른다 \cite{mcquarrie1976}.
실제로 페르미온 준위 하나를 대정준으로 풀면 상태 합은 $1+e^{-\beta(E-\mu)}$ 의 두 항이 전부다 ---
``한 상태에 $0$ 또는 $1$''이라는 셈법이 같으면 유도의 대수도, 분포의 꼴도 같아진다. 격자 삽입 자리의
점유 $0/1$ 은 페르미온과 이 셈법 구조만 공유한다 --- 배타의 근거가 파동함수의 반대칭성이 아니라
정전$\cdot$입체 반발이라는 \emph{기하}이고 입자도 전자가 아니라 Li 이온인데, 셈법이 같으므로 같은
logistic(Fermi--Dirac 형) 대수가 나온다. 진짜 양자 통계가 일하는 곳은 Chapter 2 다 --- 격자 진동의
양자(포논)는 보손이라 진동 엔트로피가 Bose--Einstein 들뜸수로 적히고, 전도 전자는 페르미온이라
electronic 엔트로피가 Fermi--Dirac 분포로 적힌다 \cite{mcquarrie1976,ashcroftmermin1976}.
\end{bgbox}

\textbf{(a$'$) 출발 --- 미시상태의 온전한 기술.} 원형의 점입자 이상화는 실제 자리에서 성립하지 않는다 ---
점유수 $n$ 만으로는 미시상태가 다 지정되지 않기 때문이다. $n=1$ 인 자리에서 이온은 한 점에 얼어붙어
있는 것이 아니라 자리 우물 안에서 진동하는 연속 자유도(위치$\cdot$운동량)를 가진다. 따라서 미시상태의
목록을 ``빈 자리($n=0$, 에너지 $0$ 기준)'' 하나와, ``점유 자리($n=1$, 우물 바닥 에너지
$\varepsilon_0$)에 내부 배치 $\Gamma$ 가 붙은 연속족''으로 넓힌다 --- 가정 1$\cdot$2 는 그대로 둔다.

\textbf{(b$'$) 연산 --- 점유수 합 $\times$ 내부 자유도 적분.} 식~\eqref{eq:sm-gc} 의 대정준 합을 이 목록
그대로 쓰면, 점유수에 대한 합(가정 1 로 두 항)과 각 점유 상태의 내부 배치에 대한 적분으로 갈라진다:
\begin{equation}
\Xi_1\;=\;\sum_{n=0,1}e^{\beta\mu n}\,z_n
\;=\;1+q(T)\,e^{-\beta(\varepsilon_0-\mu)},
\qquad
q(T)\equiv\int e^{-\beta H_\mathrm{int}(\Gamma)}\,\frac{\dd\Gamma}{h^3}
\label{eq:partfn}
\end{equation}
--- $z_n$ 은 점유수 $n$ 인 상태들의 canonical 분배함수로, $z_0=1$(빈 자리는 상태 하나),
$z_1=e^{-\beta\varepsilon_0}q(T)$(우물 바닥 인자에 내부 자유도의 적분 $q(T)$ 가 곱해진다)이다.
$q(T)$ 는 점유된 이온의 진동 등 국소 자유도의 단일 자리 분배함수다\footnote{이 $q(T)$(단일 자리 내부 분배함수)는 N6--N9 의 정규화 용량 좌표 $q{=}Q/Q_\cell$ 와 다른 양이다 --- 전자는 Part 0 의 국소 열역학량, 후자는 곡선 좌표.}. 일반적인 정의는 상태
합 $q(T)=\mathrm{Tr}\,e^{-\beta H_\mathrm{int}}$ 이고 식~\eqref{eq:partfn} 의 위상공간 적분은 그 고전
극한이다 --- 자리 우물을 3차원 조화 퍼텐셜(진동수 $\omega_0$)로 근사하면 고전 극한에서
$q=(k_BT/\hbar\omega_0)^3$ 이고, 등방 가정(세 방향이 같은 $\omega_0$)을 풀어 데카르트 축마다 다른
진동수를 허용하는 일반형으로 넓히면 양자 상태 합으로는
$q=\prod_{i=1}^{3}[2\sinh(\beta\hbar\omega_i/2)]^{-1}$(축 지표 $i=1,2,3$ 마다 진동수 $\omega_i$)로 닫힌다.
내부 분배함수 $q(T)$ 를 포함한 이 확장 역시 국소화 흡착 통계의 표준 전개이고
\cite{hill1960,fowler1939,mcquarrie1976}, $q\equiv1$(내부 자유도 없음)로 두면
식~\eqref{eq:partfn} 는 원형~\eqref{eq:sm-baresum} 으로 되돌아간다. 위 첨자 없는 $\Xi_1$ 은 내부
자유도까지 포함한 완전한 단일 자리 대정준 분배함수이고, Chapter 2 의 단일 자리 분배함수(그 문건 식
(eq:Z1))도 같은 기호 $\Xi_1$ 로 통일한다.

\textbf{(c$'$) 중간식 --- 평균 점유와 계수 $q(T)$ 의 흡수.} 점유 상태들의 확률 합이 곧 평균 점유다:
\begin{equation}
\langle n\rangle=0\cdot P_0+1\cdot P_1
=\frac{q(T)\,e^{-\beta(\varepsilon_0-\mu)}}{1+q(T)\,e^{-\beta(\varepsilon_0-\mu)}},
\label{eq:sm-occmid}
\end{equation}
그리고 식~\eqref{eq:sm-gc} 의 셋째 식으로 교차검증하면 $k_BT\,\partial\ln\Xi_1/\partial\mu$ 가 같은
결과를 준다. 여기서 적분이 남긴 계수 $q(T)$ 는 \emph{분자와 분모의 점유-상태 항 양쪽에 같은 인수로}
붙어 있다. 따라서 결과는 $q$ 와 $\varepsilon_0$ 을 따로 아는 데 의존하지 않고 한 조합에만 의존하며,
그 조합은 지수 안으로 흡수된다:
\begin{equation}
q(T)\,e^{-\beta\varepsilon_0}\;=\;e^{-\beta\tilde\varepsilon(T)},
\qquad
\tilde\varepsilon(T)\;\equiv\;\varepsilon_0-k_BT\ln q(T)
\label{eq:sm-epstilde}
\end{equation}
--- $\tilde\varepsilon$ 는 점유 자리의 \emph{자유에너지}(우물 바닥 에너지 $+$ 내부 자유도의 자유에너지
$-k_BT\ln q$)다. 자리 점유의 자유에너지 비용을 $\Delta\mu\equiv\tilde\varepsilon-\mu_\mathrm{Li}$ 로
줄여 쓰면 식~\eqref{eq:sm-occmid} 는 $e^{-\beta\Delta\mu}/(1+e^{-\beta\Delta\mu})$ 가 되어, 내부
자유도가 있어도 함수형은 원형 박스~\eqref{eq:sm-bare} 의 두-상태 logistic 그대로 보존된다.

\textbf{(d$'$) 박스 --- 점유율 $\theta$.} 분모·분자를 $e^{-\beta\Delta\mu}$ 로 나누면
\begin{equation}
\boxed{\;\theta\;\equiv\;\langle n\rangle\;=\;\frac{1}{1+e^{+\beta\,\Delta\mu}}\;,
\qquad \Delta\mu\equiv\tilde\varepsilon(T)-\mu_\mathrm{Li}\;.}
\label{eq:fermifn}
\end{equation}
\begin{verifybox}
극한 검산 --- (i) $\beta\to\infty$($T\to0$)에서 $\Delta\mu\gtrless0$ 에 따라 $\theta\to0/1$ 의 계단
($T\to0$ 에서 $\tilde\varepsilon$ 는 점유 상태의 바닥 에너지로 수렴한다 --- 고전 $q$ 면 $k_BT\ln q\to0$
이라 $\varepsilon_0$, 양자 조화 $q$ 면 영점 에너지를 더한 $\varepsilon_0+\tfrac32\hbar\omega_0$ --- 어느
쪽이든 상수라 계단 구조는 동일). (ii) $\beta\to0$($T\to\infty$)에서는 $\beta\Delta\mu\to-\ln q$ 이므로
내부 자유도가 없는 기준($q\equiv1$)에서 $\theta\to\tfrac12$(두 상태 등확률)이고, $q(T)$ 가 온도와 함께
증가하는 실제 조화 우물에서는 점유 상태의 위상공간 증가가 이겨 $\theta\to1$ 로 향한다.
(iii) $\mu_\mathrm{Li}\to\pm\infty$ 에서 $\theta\to1/0$. (iv) 확장을 끄는 극한 $q\equiv1$ 에서
$\tilde\varepsilon\to\varepsilon_0$·$\Delta\mu\to\varepsilon_0-\mu_\mathrm{Li}$ 가 되어
식~\eqref{eq:fermifn} 는 원형 박스~\eqref{eq:sm-bare} 의 $\langle n\rangle$ 을 그대로 회수한다.
★함수형 가드 --- 식~\eqref{eq:fermifn} 은 Fermi--Dirac 함수와 \emph{동형}이지만 물리량은 다르다:
공유하는 것은 ``한 자리에 입자 0 또는 1''이라는 배타 점유의 대수 구조뿐이고(가정 1), 여기의 입자는
전자가 아니라 Li 이온이며 배타성의 근거도 페르미온의 Pauli 배타 같은 양자 통계가 아니라 자리 하나에
이온 하나라는 기하다.
\end{verifybox}

한 걸음 더 --- $n\in\{0,1\}$ 라 $n^2=n$ 이므로 점유의 요동(분산)이 닫힌 꼴로 나온다:
\begin{equation}
\mathrm{var}(n)=\langle n^2\rangle-\langle n\rangle^2=\theta(1-\theta),
\qquad
\frac{\partial\theta}{\partial(\beta\mu)}=\theta(1-\theta)
\label{eq:sm-flucres}
\end{equation}
(둘째 식은 식~\eqref{eq:fermifn} 의 $\beta\mu$ 미분). \emph{점유의 $\mu$-감도 $=$ 점유 요동} --- 요동--응답
(fluctuation--response) 관계의 최소 사례이고, 본론 평형 peak 의 종 $\xi_\eq(1-\xi_\eq)$(\S\ref{sec:eqpeak})의
통계적 정체가 바로 이 분산이다.

유효 자리 자유에너지 $\tilde\varepsilon(T)$ 가 하류에서 가는 곳은 두 갈래다. 첫째, 온도를 고정하고 보면
$\tilde\varepsilon$ 는 상수 하나이고, \S\ref{sec:sm-electro} 의 전기화학 결선에서 전이 중심 $U_j$ 로
흡수된다 --- 그 소절과 본론이 쓰는 몰당 기준 에너지 $\mu^0$ 는 정확히 $\tilde\varepsilon$ 의 몰 환산
$N_A\tilde\varepsilon$ 이다. 둘째, 온도를 미분하면 내부 자유도의 자유에너지 몫
$f_\mathrm{int}=-k_BT\ln q$ 가 자리당 엔트로피
\begin{equation}
s_\mathrm{int}\;=\;-\frac{\partial f_\mathrm{int}}{\partial T}
\;=\;k_B\,\frac{\partial\,(T\ln q)}{\partial T}
\label{eq:sm-sint}
\end{equation}
를 주는데, 이것이 Chapter 2 의 vibrational 엔트로피 사슬의 단일 자리 원형이다 --- 여기 자리 하나의
세 축 진동수 $\omega_i$($i=1,2,3$)를 격자 전체의 모든 진동 모드를 헤아리는 지표 $k$ 로 넓혀, 조화 모드의
$q_k=[1-e^{-\beta\hbar\omega_k}]^{-1}$ 를 식~\eqref{eq:sm-sint} 에 넣으면 그 문건의 모드당 엔트로피
$s_k=k_B[(1+\bar n_k)\ln(1+\bar n_k)-\bar n_k\ln\bar n_k]$(Bose--Einstein 들뜸수 $\bar n_k$)가 그대로
나온다. 기준점에 관한 사소한 주의로, 위 (b$'$) 의 $[2\sinh(\beta\hbar\omega/2)]^{-1}$ 은 우물 바닥 기준
(영점 에너지를 $q$ 에 포함)이고 여기의 $[1-e^{-\beta\hbar\omega}]^{-1}$ 은 영점 에너지를
$\varepsilon_0$ 쪽에 포함시킨 기준이다 --- 두 규약의 $q$ 는 인자 $e^{-\beta\hbar\omega/2}$(영점 에너지
$\hbar\omega/2$ 의 Boltzmann 인자)만큼 다른데, 이 인자가 $T\ln q$ 에 남기는 기여는 온도 무관 상수
$-\hbar\omega/(2k_B)$ 이므로 온도 미분, 곧 엔트로피는 동일하다. 따라서 전이 중심의 온도 기울기
$\dd U_j/\dd T=\Delta S_{\rxn,j}/F$(\S\ref{sec:center})에서 진동
엔트로피가 차지하는 몫의 미시적 기원이 바로 $\tilde\varepsilon(T)$ 의 온도 의존, 곧 $q(T)$ 다.

자리 하나가 준 답 --- 점유율 $\theta$, 요동 $\theta(1-\theta)$, 유효 자리 자유에너지
$\tilde\varepsilon(T)$ --- 을 손에 쥐었으므로, 다음 소절은 같은 답을 $M$ 개의 독립 자리로 복제해 거시
점유율과 섞임 엔트로피로 넓힌다.
% ==== 블록 끝 ====
% NOTE:
% Read Coverage: AUTHOR_BRIEF.md 전문(43행) · _sections/ch1_sec02a_part0.tex 전문(L1–268, 특히 L106–219) ·
%   _sections/ch1_sec02b_part0.tex L1–70 · _sections/ch2_sec01_partition.tex L1–60 ·
%   results/V1020_STYLE_RUBRIC.md 전문(51행) · _sections/ch1_preamble.tex 전문(bgbox=\newtheorem*{bgbox}{배경}, 제목 인자 미사용 확인) ·
%   _sections/ch1_bib.tex L9–12(hill1960·fowler1939·mcquarrie1976·mckinnon1983 실재) · results/V1020_REFERENCE_LEDGER.md(ashcroftmermin1976 V1 등재·장 수준만·쪽수 금지)
% 보존 체크: P-1=(a) itemize 가정 1·가정 2(원문 문장 보존) | P-2=(a') 연속 내부 배치 Γ 붙은 미시상태 목록(원문 문장 보존) |
%   P-3=eq:partfn=(b')·eq:sm-occmid=(c')·eq:sm-epstilde=(c')·eq:fermifn=(d') 박스·eq:sm-flucres=요동 단락·eq:sm-sint=하류 단락 — 6종 전부 존재, 수식 내용 원문 그대로 |
%   P-4=(b') footnote(q(T) vs q=Q/Q_cell) 원문 그대로 | P-5=(b') 고전 (k_BT/ħω₀)³·양자 ∏_{i=1}^{3}[2sinh(βħω_i/2)]^{-1} |
%   P-6=최종 verifybox ★함수형 가드(원문 유지 + Pauli 배타/페르미온 어휘로 bgbox 와 내용 접속) | P-7=요동 단락 var(n)=θ(1-θ)=∂θ/∂(βμ)+요동--응답 문장 |
%   P-8=하류 두 갈래 단락(U_j 흡수·μ⁰=N_Aε̃·s_int=k_B∂(T ln q)/∂T·Ch2 모드 확장 s_k·영점 에너지 규약 주의 — 원문 전량 보존, (b)→(b') 지시만 갱신) |
%   P-9=(c) 원형·(c') 확장 양쪽에서 eq:sm-gc 셋째 식 교차검증 | P-10=(b) Ξ vs canonical Z 앙상블 구분 + (b') Ch2 (eq:Z1) 동일 기호 Ξ₁ 통일 문장
% 신규 라벨: eq:sm-baresum · eq:sm-baremid · eq:sm-bare (기존 라벨 6종 개명·삭제 없음; 신규 인용 키 없음 — 화이트리스트 5종만 사용)
% 자체검수: (1) 극한 — 원형 verifybox (i)β→∞ 계단·(ii)β→0→1/2·(iii)μ→±∞ + 부호 읽기(B3 충족); 확장 verifybox (i)–(iii) 원문 유지 + (iv) q≡1 ⇒ eq:sm-bare 회수 명시(B4 원형 회수 검산).
%   (2) D7 2단 분리 — 원형 박스(eq:sm-bare, 문헌 병기)와 확장 박스(eq:fermifn) 분리, 확장 동기 (a') 도입 2문장, 중간식 사슬 (b)(c)·(b')(c') 로 점프 없음(B1·B2).
%   (3) rubric — bgbox 7문장 자족(굵은 도입 시작·지시대명사 無·Ch2 forward 다리 포함), 자기 개정 서술 0·명령형 0·코드 언급 0, 용어 병기(페르미온(fermion)·보손(boson)·Fermi--Dirac·Bose--Einstein), 절 도입·마무리 다리 존재(A3).
